
THE FREE CHURCH’S SECOND ANNUAL MEETING

—

(Continued from the preceding No.)

\subsection{Friday Morning.}

Prayer meeting from 9:00 to 9:30, led by Pastor B. T. Petersen.

The minutes of the previous meeting were read and approved.

A communication from the congregation in Seattle, Wash., through Pastor B. A. Borrevik, was read.

They then proceeded to consider the draft of a new Congregational Order.

“Congregational Order
for
— — — Congregation — — —

Article 1.

The name of this Congregation is — — —

Art. 2.
Confession.

§ 1. This Congregation believes and confesses that Holy Scripture, the canonical books of the Old and New Testaments, is the Word of God, revealed for the salvation and blessedness of mankind, and is the sole source and rule for faith, doctrine, and life in the Congregation.

§ 2. This Congregation confesses the ancient ecumenical Symbols, Luther’s Small Catechism, and the Unaltered Augsburg Confession.

Art. 3.
The Congregation.

§ 1. The Congregation is the Body of Christ. Col. 1:24. It is the city on the hill, and cannot be hidden (Matt. 5:14), the Zion from which God shines forth. Ps. 50:2.

As such, its members must, through the use of the means of grace and by faith, be united with the Congregation’s Head and Lord, Jesus Christ.

§ 2. Every member of the Congregation ought, according to his measure of faith, to use the gift of grace given by God to each for the edification of the Congregation and the extension of God’s Kingdom, according to the order and direction given by God’s Word, and also, according to his ability, to share in the expenses of the Congregation.”

The paragraphs quoted above were adopted without debate.

Art. 4.

\subsubsection{Membership and Voting Rights.}

§ 1. Adults who are received into the Congregation must:

1. have been baptized with Christian Baptism,

2. show in confession and life that they seek first the Kingdom of God and his righteousness,

3. assent to the Constitution of this Congregation.”

Mr. Hofstad wished to know whether no distinction was to be made between women and men.

The chairman explained that this question would be dealt with under a later point.

Prof. Sverdrup thought that in item 2 a peculiar passage of Scripture had been inserted—merely a general exhortation, which set no definite condition for admission into the Congregation.

Pastor Aas: The committee was not agreed on this point. The wording is unfortunate. I have not been able to go along with the point. The passage of Scripture contains only a general exhortation and is not a sufficiently clear expression of what ought to be required as a condition for admission into the Congregation.

It can be understood in so many ways. I therefore have a minority recommendation, which reads as follows: “in life and confession testify that through faith they are united with the Head of the Congregation, Jesus Christ.”

I think this is more consistent with what has been proposed concerning confirmands—that they must have come to conscious faith.

If this is required of confirmands, then it seems to me that one ought also to require at least so much of adults who wish to enter the Congregation.

Delegate Blomlie thought there was an inconsistency between this paragraph and a later one. It is not laid down here as a condition for admission into the Congregation that one must have been confirmed, while it is made a condition for obtaining the right to vote.

Pastor Winther believed it was of great importance that this point be thoroughly discussed. From one point of view, it may not matter so greatly what is adopted here; for the draft is only a model, not a binding law for the congregations. Yet it is good to discuss the matter.

The exceedingly difficult point is the question of the admission of adults into the Congregation. The minority recommendation seemed to require a conscious life with the Head, Jesus Christ. To this belongs also, in my opinion, participation in the Lord’s Supper; but how is all this to take place with a person outside the Congregation?

I do not believe that a person can come to a conscious life with the Head, Jesus Christ, outside the Congregation.

We have also previously had good conditions for admission into the Congregation; but it is one thing to have them on paper, another to put them into practice. I am, however, not satisfied with the old order; I therefore rejoice over every step in the direction of reform. Yet I do not believe it is so imperative to include these definite rules with a requirement for these definite marks. What matters, rather, is the inward spiritual work. If vigorous spiritual work is done in the Congregation, people will as a rule not apply for admission before they also have the proper qualifications for it.

If spiritual need is present, then I do not think one ought to require that they have come to a conscious faith. There is also something in the Christian life called longing faith; therefore I cannot go along with the minority recommendation.

Jonas Helland: It may perhaps seem strange that we could not agree in the committee. We could, however, agree well enough concerning the wording of the minority recommendation, but not concerning its interpretation. We did not believe that the penitent and those who are seeking ought to be excluded from the Congregation, even though they lacked confident, conscious faith.

Faith has its development; longing and seeking are signs of faith beginning. There are those who never seem to come to this confident, conscious faith, but who nevertheless show in their lives that they walk with the Lord. The Savior himself calls those blessed who hunger and thirst for righteousness. Why, then, should we not give such persons access to the Congregation?

Theod. Johnson: The previous speaker has in substance presented what I wished to say. It is true that many are seeking for a long time before they come to conscious faith; so it was with me. I think it is harsh to exclude those who are seeking.

I would therefore express it in this way: “who show that they are truly seeking.”

Pastor Arevig: I think a different requirement is set forth in III, 1 and IV, 2. Those who are admitted ought to stand approximately on the same spiritual level as those who admit them; otherwise there will be no true union. I therefore favor the minority recommendation, since it refers back to III, 1.

Pastor Midthun: This point comes very near the same matter that we discussed under “guiding principles.” I will say here also that one is in danger of going too far. I fear that this will lead to carnal zeal and spiritual domineering. I think III, 1 and IV, 2 speak of two different things; the former point speaks of how one becomes a child of God, the latter of what conditions are to be set for admission into the Congregation. I am more inclined toward the majority recommendation.

Pastor Kjeldas thought the point was strict enough without any amendment. He wished especially to call attention to the difficulties on the mission field. How can one there demand conscious faith as a condition for admission into the Congregation? And how can the missionary pastor come to know the spiritual condition of the individual? He is most often unacquainted with them and comes there only seldom—who, then, is to testify to the genuineness of the people’s Christianity?

Moreover, one can so easily be mistaken and judge wrongly; I have myself experienced this. I am afraid one may get onto a slippery slope, become like the “Free-Free,” and pass over both repentance and conversion.

Pastor Birkeland: I have exceedingly little faith in written rules. As the members of the Congregation are, so the Congregation is—not as the rules are.

In this connection I shall quote a statement by Prof. Böckman; I do so all the more readily because it comes from a quarter from which one would least expect it.

Prof. Böckman said that all congregations in this country, so far as he knew them, were based upon a lie. They begin with a lie by admitting anyone and everyone without asking whether they are believers.

Nevertheless, it is good to have an expression of how members of the Congregation are to be admitted. The principal matter, however, is not to guard oneself against unbelief, but that all believers be included. The Congregation cannot afford to leave any believing soul outside. Christ’s Body and Blood cannot be cut apart; and when we exclude those who belong to Jesus’ Body, what are we doing other than cutting it apart?

The expressions in the point are too indefinite. Everyone will say that they seek first the Kingdom of God and his righteousness, but each understands it in his own way.

But then comes the question: Have we the right to require of those who seek admission that they must be able to answer with a definite Yes that through faith they are united with Christ? The principal matter in the confession is not that one can stand up and say: I am converted; I believe—it is all well and good, and I rejoice over everyone who can confess it; neither am I so terribly afraid of the “Free-Free”—but the principal matter is not what one can confess about oneself; the question is: What do you think of Christ? One must be able to say something about Jesus, confess something about him. There is altogether too little testimony concerning Jesus in our time—even in public preaching.

If there is to be a condition, then let it be this: What can you say about Jesus? What may a person not confess about himself when hard pressed!

No, let them confess something about Jesus—even if no more than the old woman who said: I cannot say much about him; but this I know, that I love him so that I would gladly die for him.

Prof. Sverdrup: One can also speak falsely about Jesus. One may adopt certain insipid phrases about Jesus and mean very little by them. So neither is that any sure sign of true life in God.

The difficulty is not admission into the Congregation, but the difficulties must be considered in connection with the Congregation itself. Who is to admit? What is the Congregation like?

When the Congregation is worldly, and the pastor is worldly, they naturally admit such as are like themselves. If, on the other hand, one had an assembly of believers and a believing pastor, there would not be such great difficulties with admission into the Congregation.

As conditions are in most places, it is simply impossible to bring matters into proper order. An abuse is prevalent. Many enter the Congregation in order to obtain a security with respect to eternity; others wish to obtain protection against attacks upon sins and vices within the Congregation. When, for example, a pastor takes a stand against youthful excesses, people are brought into the Congregation in order to drive the pastor away.

It is therefore right to be quite exacting about admission into the Congregation; above all, one must admit man by man, person by person. Admission in this connection means nothing other than that one gives those who are admitted an opportunity to participate in governing the Congregation, not that one prescribes to them how they are to be Christians or gives them access to hearing God’s Word.

The stricter one is, the better. One thereby prevents no one from salvation and blessedness.

A person can be saved even though his name is not in the congregational register.

From God’s Word we see that those who were baptized and saved were also added to the Congregation. If, now, a congregation wished to draw a boundary and say: only when you have come to conscious faith can you participate in governing the Congregation, it surely has the right to do so.

I would prefer to word the point differently, so that it might agree with VII, 4, where it says that the Congregational Council shall “investigate whether those who request admission into the Congregation possess the necessary qualifications for it.” There ought to be some testimony, and preferably the testimony of those who have examined the person concerned.

Pastor Broen: The principal matter at this meeting is not to establish something definite, but that a great many thoughts may be brought forth and shed light upon the matter; for that is needed. We need counsel and direction concerning how one ought to proceed in admitting new members.

Something new is pressing its way forward. And it is probably now as before, that the new wine will burst the old wineskin. The believers in the Congregation see that things are wrong as they are. And that is true. There must be some testimony concerning what Christ has been able to accomplish in them.

Suppose now that the Congregation has elected believing deacons. Applications for admission are presented to them, and they are to act according to their new light; then the difficulties arise, for there comes to be such a striking contrast between the new and the old. And there will be severe conflicts when one begins to put the new into practice.

I have never liked the committee’s recommendation; I do not think it is much better than what we have had before.

If one were without more ado to admit seekers, one would get a great many wavering, indefinite persons who would not be added to the working force against evil, but would instead occupy an indefinite neutral position. If greater conflicts arise, they nevertheless tend to place themselves on the side of the world.

\subsection{Friday Afternoon}

Devotion by Pastor M. A. Pedersen of St. Paul.

The Credentials Committee reported 16 new members of the meeting.

It was resolved to give Rev. Fenton of St. Paul an opportunity to speak for ten minutes concerning secret societies.

It was resolved to take an offering for the Deaconess Home at the evening service.

Pastor O. L. Jordvik proposed that it should become the practice among the pastors in the Free Church to divide travel expenses equally.

Prof. Nybahl thought that this could not simply be done by a hasty resolution, and therefore proposed that a committee of three members be appointed to take the matter under consideration and later report to the meeting.

The latter proposal was adopted.

The following were elected to the committee: Pastor O. L. Jordvik, Pastor M. A. Pedersen, and Pastor J. H. Brønos.

Pastors K. O. Lundeberg, J. Tollefson, and S. M. Stenby were admitted as advisory members of the meeting.

At the proposal of Pastor Broen, a collection was taken for the benefit of a sick man from Stordal Congregation. The man was poor and had recently come to the city to seek medical help.

They then proceeded to the order of the day, which was the continued consideration of the draft of the Congregational Order.

Pastor P. Nielsen: There seems to be great anxiety lest we should come to admit too many into the Congregation. My anxiety has not precisely been that; what concerns me is the Congregation itself, what it is like. If the Congregation has its proper constitution, if it is led by the Spirit of God, then it will surely also be quite careful in admitting new members.

It may be that the committee has not been able to express itself rightly concerning the Congregation; but it is one thing to express ideally what it ought to be, another to realize it. What I, for my part, had in view was this: that the Congregation ought to place itself on the side of the Spirit of God and not in opposition to the Spirit or his work.

The Spirit does not work only within a congregational organization, but everywhere God’s Word sounds forth. If the Holy Spirit has begun his work in a human heart, that person will also be drawn to those who belong to the Christian Congregation; I have experienced it myself.

If one now asks the person seeking admission: What is your confession? and he answers: I cannot give any definite confession; but I seek Jesus—would it then be right to deny such a person admission into the Congregation? I do not believe so. Perhaps such a person is every bit as sincere as the one who can give a confident testimony.

One can also force that kind of confession and get people to bring forth with their mouths the desired testimony without possessing in the heart what corresponds to it.

I know that there is danger here also. Carnal congregations will misuse this point, just as they will misuse any other point one adopts.

It is, however, not human beings but the Spirit of God who adds to the Congregation. The Lord adds to the Congregation those who are being saved. But cannot these people wait and remain outside until they come to a greater measure of faith? one says.

(Continued on page 4.)

THE FREE CHURCH’S SECOND ANNUAL MEETING

—

(Continued from page 3.)

We must seek to stand in agreement with the Spirit of God and his work. If the Spirit adds them, we ought not set ourselves against it.

I have seen how they have proceeded where conscious faith is required. They have produced a kind of spiritual excitement and thereby forced forth the desired confession, but in a manner that has become harmful to the work of God’s Kingdom.

Ole Pedersen: It has become clear to me, from what I heard yesterday and today, that people fear we have now become so Christian that we are going to become too strict in admitting new members. I do not believe, however, that the minority has become quite so Christian yet. One may be sure that wherever the worldly-minded are in the majority, they will continue to be quite lax. Until there is awakening and quickening in the congregations as they now are, there will be no danger in that respect.

Until now people have not asked about Christianity, but rather about money. The situation is now such that I can see no way out. So long as all the work is left to the pastors alone, little can be done; for as a rule they are overloaded and have too large a field of labor. I have often thought of what Skrefsrud said about how they do it in Santalistan, that there they go out two by two to speak with people. Think if we did the same? Can we?

Pastor Quambeck: The question before us is to obtain a rule in accordance with God’s Word, whether the congregations are sleeping or not.

We are surely agreed that all who are to be added to the Congregation ought to belong to the Lord. I also believe that they ought to show it in confession and life.

It was resolved that a vote should be taken.

It was proposed that the point be referred back to the committee. Defeated.

The point—the majority recommendation—was adopted.

“§ 2. All children become members of the Congregation through Baptism.”

Ole Pedersen: It is a well-known fact that pastors baptize children of parents who stand outside the Congregation; is this to be understood to mean that these children also become members of the Congregation through Baptism? I have never liked pastors having various little jobs outside the Congregation. Did the first Christians do anything of the kind? We act thoughtlessly.

For a long time I thought this was something the pastors did on their own. But then I asked our own pastor about it. He said that he could not do such a thing at all except as the servant of the Congregation.

Are children of outsiders to become members of the Congregation through Baptism?

Pastor Helland wished to ask the same question.

It would be clear and intelligible if there were added: “of parents who belong to the Congregation.”

As a pastor, he wished to receive guidance in this matter. City pastors are exposed here to distress of conscience, more so than country pastors. He hoped there would be a clear and definite understanding of this difficult question.

Are we to baptize children of godless parents, and are such children then to belong to the Congregation?

Pastor D. Paulson: This is an old story. The church bodies have never succeeded in regulating the matter. It applies to weddings and funerals just as much as to Baptism. If we do not perform these things, others will.

God’s Word says: Let the little children come to me and do not hinder them, for the Kingdom of God belongs to such as these.

Now I know that the child becomes a child of God through Baptism; is it then right for me to deny it the Kingdom of God because of the parents? I am so situated that I even sometimes receive requests from Germans to baptize their children, and I do it; but they do not thereby become members of our Norwegian congregations.

Pastor Winther thought it was painful to hear people call baptizing children “jobbing”—a rather ugly insinuation. In this matter I stand approximately at the same point as Pastor Paulson. I have not been able to say no to the Baptism of little children.

Pastor P. Nielsen: The point is clear enough. I will only ask: Are the children members of the Congregation before they are baptized? Do they become members of the Congregation when they are baptized and brought to Christ?

The question here is not one of practice, but of God’s Word, and it says: Go therefore and make all nations my disciples, baptizing them in the name of the Father and of the Son and of the Holy Spirit, and teaching them to keep all that I have commanded you. Here we are commanded not only to baptize, but also to teach. Shall one then baptize children indiscriminately without also seeing to their instruction in God’s Word?

I believe there is some limitation upon the use of Baptism; there are cases in which the Congregation must say no.

Can these parents who stand outside not baptize their children themselves? Why should the pastor assume a responsibility that he need not take upon himself? Let these people baptize their own children if they wish to have them baptized. Or do we Lutherans perhaps teach that it is no valid Baptism unless the pastor performs it?

We do not teach that. We ourselves are to blame that many people wish to make use of pastor and Congregation for certain “jobs,” while at the same time in their hearts they despise both pastor and Congregation.

People say: If we do not do it, then others will—is that any argument for doing what is wrong? We must give some thought to the responsibility involved in the use of the means of grace.

If we determine that children “of parents who belong to the Congregation” become through Baptism members of the same, is that then to mean that children of other parents stand differently? Does not Baptism work the same in both cases? Does it work in two different directions? Baptism always works in the same way, whether the parents belong to the Congregation or not.

Pastor Chr. Jirehus held the same view as Nielsen. If the Congregation baptizes children, then it also has responsibility that they be instructed in Christian doctrine for children and in God’s Word. It ought also to be regarded in this way, that children baptized in the midst of the Congregation are members of it; for it is altogether easier to keep oversight of those who belong to the Congregation than of those who stand outside. And it ought to be understood that when the Congregation baptizes children, it will also have the right to bring them up.

Baard Andersen was against drawing the bonds too tight. One can so easily place hindrances in the way of the work of God’s Spirit in the Congregation, hindering the little ones from coming to Christ. To illustrate what he meant, he told of a case from his own congregation.

There was a married couple who did not belong to the Congregation. The man associated with the sects; but the mother adhered to the Lutherans. She also wished to have the children baptized. They were baptized. As they grew older, they were also sent to the Lutheran religious school. One of the children, now a young man, stands today as a respected member of the Congregation and is a great blessing to it. Would this have been the case if the Congregation had denied these children Baptism?

I am not in favor of setting bounds to the activity of pastors; but I believe it is best to let the pastor act according to his conviction.

Let us proceed cautiously and not go forward too quickly or too violently; for if one takes things too abruptly and too forcefully, conflict and division will follow.

Pastor E. Hansen: There is no doubt that this question of baptizing children of parents who stand outside the Congregation becomes a matter of conscience. The worst thing is that one is tempted to try to make God’s Word fit what is wrong.

I have read about the Catholics’ missionary work, how they steal children, baptize them, confirm them, and make them Catholics for life. Are we to practice something similar?

I am in favor of the children becoming members of the Congregation through Baptism, and of their being reminded of it when they grow up.

Jonas Helland: What we in the committee meant was that when the pastor baptizes children on behalf of the Congregation, they become members of the same. What he does on his own account is his affair. For that matter, I doubt whether he has any call at all to baptize without the consent of the Congregation.

Pastor Aas: It is not stated in the point that one is not to baptize children of parents who stand outside the Congregation. What we thought in the committee was this: that one had no right to baptize children over whom one could afterward exercise no oversight. The Congregation ought not baptize children of parents who neither will themselves bring up their children in Christianity nor permit the Congregation to do so. That others baptize them if we do not is something we cannot help.

It sounds ugly to call it “jobbing”; but the worst of it is that it is true.

Theo. Johnson: The point is unclear. I would like it better if it had said that all children become “born again,” or “members of the invisible Church.”

There is also a difference among those who stand outside. If I were a pastor, I would have the parents sign a declaration in which they promised to bring up their children in the Christian faith.

It was resolved to cut off the debate and proceed to a vote.

Point 2 was adopted.

“§ 3. Entitled to vote in this Congregation are all members who are over 21 years of age, confirmed, and not under congregational discipline.”

Pastors Thomasberg and Broen each had an amendment to the point. The former’s received no second. Broen wished to insert the words “who have access to the Lord’s Supper” instead of “who are over 21 years of age.” His amendment was seconded.

Pastor Yderstad: I do not like the word “all” members; I therefore propose that the word “male” be added.

I am still so conservative that I do not believe it is right to give women the vote. I believe it is good for us to take a little time and consider what we are doing before we proceed too hastily. The Free Church is surely the most progressive among the Norwegian-Lutheran church bodies in this country, and I am glad of it; but we are in danger of becoming too progressive. And it must not be taken amiss if someone tries to hold back. I, for one, will try to do so. Let there be a little calm and steady progress among us. This question will surely awaken conflict, and we do well to think of that also. I propose that we conclude consideration of this draft for this year.

The proposal failed.

Mr. Hofstad knew of congregations that had introduced voting rights for women. No unpleasantness had arisen on that account; rather, it had worked for blessing.

Prof. Oftedal: I am one of the conservatives and can well understand the misgivings some have.

We must make ourselves familiar with the fact that reforms are of little use unless they are borne by personal conviction.

At the same time, we must remember that our provisions are binding upon no one, but are only counsel or guidance, and therefore ought not to be feared.

I am strongly opposed to the amendments that have been made. They concern congregational discipline and not qualifications for voting. “Access to the Lord’s Supper” is a doubtful expression and altogether too indefinite. There ought to be an age limit. Whether 21 years is the most suitable may be a matter on which opinions differ. Yet this is fundamentally a conservative point, by no means so excessively progressive.

It was resolved to depart from the order of the day and hear the report of the Board of Trustees for Augsburg Seminary.

Prof. Oftedal then read a brief report, from which it appeared that, following the Supreme Court decision, Augsburg Seminary is hopefully secured for the Free Church. Augsburg’s friends will rejoice over this and thank God and will not reproach their opponents.

The lawsuit has cost Augsburg about 5,000 dollars.

The school needs larger classrooms and more religious instruction in English.

The following were elected to the committee on the Board of Trustees’ report: Pastor A. Houkom, Pastor J. H. Brønøs, Pastor G. Nilsen, J. Strandnes, and Ole Pedersen.

In the evening Pastors J. Tollefsen and S. M. Tollefsen preached in Trinity Church, and an offering was taken for the Deaconess Home.

Pastors Østgaulen and Ytrehus preached in St. Olaf’s Church, Broen and Dahle in St. Luke’s.

(To be continued.) 

%[1]: https://snl.no/Den_norske_Santalmisjon?utm_source=chatgpt.com "Den norske Santalmisjon – Store norske leksikon"
